\chapter{Fresh C14.4 completion when $\kappa$ is small and both Bob source fractions are units}
\label{component:p2edb55cd1f69}
\noindent\textbf{Component }\path{p2edb55cd1f69}\par
{\small\noindent Source: \path{NEW/c1/c144_finite_matched_rederived.tex}\par}
{\small\noindent SHA256: \path{2edb55cd1f6919d774d82054aee53db4ff6b63b6444cb535185ac06146cb308a}\par}
{\small\noindent Inventory use: FM1--FM22 | FM2–FM3 actual source formulas only; unit-fraction asymptotic scope not imported | Fresh full-parent proof FM1-FM22, candidate awaiting independent analytical review\par}
\medskip
\section{Fresh C14.4 completion when $\kappa$ is small and both Bob source fractions are units}
\label{p2edb55cd1f69:sec:new-c144-finite-matched}

\paragraph{Scope and status.}
This is a fresh candidate theorem, pending the focused algebra receipt
and independent analytical review. It does not assert completion of
the source endpoints $\nu\in\operatorname{Inf}$ or
$\ell\in\operatorname{Inf}$, or of the unit-$\kappa$ source.
All original physical C14.4 source conditions, strict gaps,
nonnegative event coefficients, total-scale conditions and prior
complement intersections are retained. No earlier missing R21 proof
body is invoked.

\paragraph{Theorem.}
Let the original native C14.4 source have
\[
 S=\sigma+L>0,\qquad x=\kappa S,\qquad
 \nu=\sigma/S,\quad\ell=L/S,\quad\nu+\ell=1,
 \qquad S,\kappa\in\operatorname{Inf}_{>0},
 \tag{FM1}
\]
with $\nu,\ell$ positive bounded units. Then the complete original
parent is a finite SOS in $K_c\otimes_{\mathbb R}K_t$ for every
independent projective test and every real or complex state.
The argument below supplies the missing finite matched-Bob chart;
the other tests use the surviving explicitly scoped SE theorem.

\paragraph{Exact opening-safe source dictionary.}
Write $\theta=1-h$, $0<r,h<1$, $e=\kappa h/\nu$, $E=1-e$,
and set
\[
\begin{split}
 J&=1+xh-\nu S(1-r),\quad
 M=1+xr-\kappa\theta/\ell,\quad Q_s=1-\ell Sr,\\
 P_0&=1+S\{\nu\ell+2\nu\kappa h+\kappa\ell+
                    \kappa^2h\theta\}
             +S^2\{\nu r\kappa\ell+\nu r\kappa^2h\}.
\end{split}\tag{FM2}
\]
These are the actual native source quantities, with $P=SP_0$.
The converted P1 coefficients, without applying P1's different
unit-$L$ asymptotic hypotheses, are exactly
\[
\begin{split}
 c_0=SC(1-r)&=\frac{(1+x)(1+\nu Sr)E}{2P_0},\\
 a_0=SAr&=\frac{(1+xh)E+(1-r)(\ell S+xh)}{2P_0},\\
 d_0=D/S&=\frac{(1+x)(1+\nu Sr)-JP_0}{2JSP_0},\\
 a_1=\alpha/S&=
 \frac{(1+\nu Sr)\{\ell(1+xr)e+\ell^2ME\}}{2Q_sP_0},\\
 b_1=\beta\theta&=
 \frac{(1+x)(1+\nu Sr)(\ell+\kappa h)-\theta P_0}{2P_0},\\
 \mu=m/\kappa^2&=
 \frac{(1+x)(1+\nu Sr)h}{2M\ell\nu P_0}.
\end{split}\tag{FM3}
\]
The numerator in the $d_0$ formula has a factor $S$, so its displayed
$S$ denominator cancels in the source polynomial ring. Thus all six
coefficients belong to the bounded source algebra generated by
$S,\kappa,h,r,\nu,\ell$ and inverses of $J,M,Q_s,P_0,\nu,\ell$.
None of these inverses is of an opening $r,1-r,\theta,D$ or $\beta$.
The original source gives
\[
\begin{gathered}
 J,M,Q_s,P_0,E\sim1,\qquad c_0,a_0\sim\tfrac12,\qquad
 d_0\sim\nu^2/2,\quad a_1\sim\ell^2/2,\\
 b_1-(h-\nu)/2\in\operatorname{Inf},\qquad
 \mu-h/(2\nu\ell)\in\operatorname{Inf},\qquad
 \bar h\ge\bar\nu>0.
\end{gathered}\tag{FM4}
\]
The last inequality follows from $b_1=\beta\theta\ge0$.
Hence $\mu$ is a positive bounded unit. These estimates are uniform
through $r,1-r,\theta\to0$; no comparison with a testing quantity
has been used. In this sector $D/S$ is a positive unit, so $D=0$
has empty intersection; $\beta=0$ is allowed.

\paragraph{Exact full parent and unrestricted normals.}
On the finite Bob chart write $D_u=1+u^2$, $D_v=1+v^2$.
The complete quarter-parent is
\[
\begin{split}
 F/4={}&m(vW-X+uvY-uZ)^2+C D_v(uW-Y)^2+A D_uD_vY^2\\
 &+(D+\tfrac12)(uW-Y+uvX-vZ)^2
 +(\alpha+\tfrac12)D_u(vY-Z)^2\\
 &+(\beta+\tfrac12)D_v(X+uZ)^2-\tfrac12D_uD_vZ^2.
\end{split}\tag{FM5}
\]
No state restriction is made. Put
\[
 N_x=X+huZ,\qquad N_y=\frac{Y-ruW}{\sqrt{Sr(1-r)}}.
 \tag{FM6}
\]
The source square root is nonzero even when its size is arbitrarily
small. It is used only to define a permitted tensor root; every
coefficient below is bounded without its inverse.

After the surviving unmatched-residue and adjacent-$C,A$ exits it is
enough to take $u$ infinitesimal and $v$ a bounded positive unit with
the same residue as $b=\sqrt{\ell/\nu}$. Set
\[
\begin{gathered}
 a=u^2,\quad j=\sqrt{\nu\ell},\quad T=\kappa W,\quad U=uZ,\\
 P=uW/\sqrt S,\quad Q=\sqrt S Z,\quad N=T-jU,\quad G=P-jQ,\\
 \delta=v-b,\quad B=\delta Q,\quad B_u=\delta U.
\end{gathered}
 \tag{FM7}
\]
Every state coordinate remains arbitrary, including zero and
unbounded coordinates. The rows satisfy the exact MB identities
\[
\begin{gathered}
 R_Z=(a-x)Z=(\kappa\sqrt S/j)G-(u/j)N,\\
 (a-x)\delta Z^2=(\kappa/j)GB-(1/j)NB_u,\\
 \delta^2N^2+2j\kappa\delta GB+2j^2\kappa B^2-j^2B_u^2
                  =(\delta T)^2.
\end{gathered}\tag{FM8}
\]
These are polynomial identities, not ordered-ring inequalities.

\paragraph{The full graph scalar.}
Let $F_0$ denote FM5 with $N_x=N_y=0$; it is used only as the
constant coefficient of the full normal expansion. Define bounded
source/testing polynomials
\[
\begin{split}
 A_T={}&\mu v^2(1+ra)^2,\\
 A_P={}&(1+v^2)\{c_0(1-r)+a_0r(1+a)\}\\
       &+S\{(D+\tfrac12)(1-r)^2+
                   (\alpha+\tfrac12)(1+a)r^2v^2\},\\
 A_Q={}&d_0v^2+a_1,\\
 A_U={}&m\theta^2+\alpha+b_1\theta+\theta^2/2\\
       &+v^2\{2Dh+h-\tfrac12+b_1\theta+\theta^2/2\}
              +(D+\tfrac12)h^2v^2a,\\
 K={}&m\theta(1+ra)+(D+\tfrac12)(1-r)(1+ha)
                         +(\alpha+\tfrac12)(1+a)r,\\
 H={}&vK-\kappa jA_T,\quad R_U=A_U-j^2A_T,\\
 g={}&jA_P-H,\quad R_Q=A_Q+j^2A_P-2jH,\\
 J_Z={}&SR_Q+aR_U.
\end{split}\tag{FM9}
\]
Direct expansion of every term of FM5 gives
\[
 F_0=A_TN^2+A_PG^2+2gGQ+J_ZZ^2.
 \tag{FM10}
\]
For example the only relation used to obtain FM10 from the
five monomials $T^2,U^2,P^2,Q^2,PQ$ is $TU=\kappa PQ$.

The actual source stationarity identities, also verified directly
from FM2--FM3, are
\[
 g(x,b)=0,\qquad
 J_Z(x,b)=\partial_aJ_Z(x,b)=\partial_vJ_Z(x,b)=0.
 \tag{FM11}
\]
These are coefficient-field identities at the actual source
contact; they impose no equation on independent tests.

\paragraph{Finite polynomial division.}
All following quotients mean monic polynomial division in formal
$a$ or $v$, performed before evaluation in the tensor ring. No
mixed inverse is introduced. There are bounded coefficient
polynomials $g_a,g_v,H_{20},H_{11},H_{02}$ defined exactly by
\[
\begin{split}
 g(a,v)&=(a-x)g_a(a,v)+(v-b)g_v(v),\\
 J_Z(a,v)&=(a-x)^2H_{20}(a,v)+(a-x)(v-b)H_{11}(v)
                                      +S(v-b)^2H_{02}(v).
\end{split}\tag{FM12}
\]
For clarity, first divide $g(a,v)-g(x,v)$ by $a-x$ and then
$g(x,v)$ by $v-b$. For the second line, first subtract the
constant and linear Taylor terms in $a-x$ and divide the remainder
by $(a-x)^2$. Divide $\partial_aJ_Z(x,v)$ by $v-b$ for $H_{11}$.
Finally divide the bounded polynomial
$R_Q(x,v)+\kappa R_U(x,v)$ by $(v-b)^2$ for $H_{02}$.
FM11 makes every division exact. In particular $H_{02}$ is
bounded without dividing a bounded error by $S$.

Substituting FM8 into FM10 now yields the exact scalar identity
\[
\begin{split}
 F_0={}&A_TN^2+A_PG^2+2g_vGB+H_{02}B^2
          +H_{11}\{(\kappa/j)GB-(1/j)NB_u\}\\
       &+H_{20}R_Z^2+2\sqrt S\,g_aGR_Z.
\end{split}\tag{FM13}
\]
Thus every coefficient is bounded and all remaining radial
contributions are infinitesimal matrix entries on $(N,G)$.

Here and below a displayed residue formula uses the real residues
of its bounded source parameters. FM3 gives explicitly
\[
 \bar A_T=\frac h{2\nu^2},\quad
 \bar A_P=\frac1{2\nu},\quad
 \bar g_v=\frac{\ell-\nu}{2},\quad
 \bar H_{02}=\frac\nu2,\quad \bar H_{11}=-\theta j.
 \tag{FM14}
\]
For example at $S=\kappa=0$ and $a=0$ the three relevant
polynomials are
$g=j(1+v^2)/2-v/2$,
$R_Q=\nu(v-b)^2/2$, and
$R_U=\theta(\ell-\nu v^2)/2$.
These establish FM14 by the finite divisions just specified.

\paragraph{Both complete first-normal terms.}
Write the full expansion of FM5 as
\[
 F/4=F_0+H_xN_x^2+H_yN_y^2+2H_{xy}N_xN_y
                  +2N_xL_x+2N_yL_y+\beta D_vN_x^2.
 \tag{FM15}
\]
The exact bounded diagonal and cross coefficients are
\[
\begin{split}
 H_x={}&m+(D+\tfrac12)u^2v^2+D_v/2,\\
 H_y={}&r c_0D_v+(1-r)a_0D_uD_v\\
       &+Sr(1-r)\{mu^2v^2+D+\tfrac12
                            +(\alpha+\tfrac12)D_uv^2\},\\
 H_{xy}={}&-uv\sqrt{Sr(1-r)}(m+D+\tfrac12).
\end{split}\tag{FM16}
\]
In particular $\bar H_x=\bar H_y=1/(2\nu)$ and
$\bar H_{xy}=0$. The last term of FM15 is retained whole;
it allows $\beta$ to be arbitrarily large or zero.

Define
\[
\begin{split}
 L_U={}&m\theta-(mj/\kappa)v(1+ra)
 -(D+\tfrac12)v^2(1+ha)+(b_1+\theta/2)(1+v^2)\\
       &+(D+\tfrac12)(1-r)vjS,\\
 Y_P={}&(1+v^2)\{-c_0+a_0(1+a)\}
       +S\{mv^2(1+ra)-(D+\tfrac12)(1-r)
                        +(\alpha+\tfrac12)(1+a)rv^2\},\\
 Y_Q={}&-m\theta a+(D+\tfrac12)(1+ha)
                                  -(\alpha+\tfrac12)(1+a),\qquad
 Y_C=jY_P+vY_Q.
\end{split}\tag{FM17}
\]
The expression $mj/\kappa=\kappa\mu j$ is bounded and
infinitesimal. The complete original normal terms are exactly
\[
\begin{split}
 L_x={}&-\kappa\mu v(1+ra)N
       +(D+\tfrac12)(1-r)vu\sqrt S\,G+L_UU,\\
 L_y={}&\sqrt{r(1-r)}\{Y_PG+Y_CQ\}.
\end{split}\tag{FM18}
\]
The actual source gives $L_U(x,b)=Y_C(x,b)=0$. Perform the same
finite monic divisions to write
\[
 L_U=(a-x)l_a+(v-b)l_v,\qquad
 Y_C=(a-x)y_a+(v-b)y_v.
 \tag{FM19}
\]
Then $L_UU=u l_aR_Z+l_vB_u$ and
$Y_CQ=\sqrt S\,y_aR_Z+y_vB$. All their coefficients are
bounded. Moreover $\bar l_v=-j$, while $Y_P$ and $y_v$ have
zero residue. This follows directly from the source-limit
identities
$L_U=(\ell-\nu v^2)/2$ at $a=0$, and
$Y_P=Y_Q=0$ at $a=0$.
Consequently $L_x+jB_u$ is an infinitesimal bounded combination
of $(N,G,B_u)$ and $L_y$ is an infinitesimal bounded combination
of $(N,G,B)$. This accounts for both unrestricted first-normal
terms, including all $\beta\theta$ terms and the independent
$r,1-r,\theta$ openings.

\paragraph{A strict residue for the full six-row Gram.}
Use FM8 with a fixed real $c=4$ to introduce a $cB_u^2$ pivot
at infinitesimal Gram cost. More explicitly subtract
\[
 (c/j^2)\delta^2N^2+(2c\kappa/j)\delta GB+2c\kappa B^2
 \tag{FM20}
\]
from the bounded part of FM15, add $cB_u^2$ to that part, and
retain the exact positive square $(\sqrt c\,\delta T/j)^2$.
Their sum is unchanged by FM8. The residue on
$(N_x,N_y,N,G,B,B_u)$ is therefore
\[
 \frac{N_x^2+N_y^2}{2\nu}+\frac h{2\nu^2}N^2
 +\frac{G^2}{2\nu}+(\ell-\nu)GB+\frac\nu2B^2
       +\theta NB_u-2jN_xB_u+4B_u^2.
 \tag{FM21}
\]
The $(G,B)$ block has determinant $\nu\ell>0$ and positive
diagonal. The remaining block has positive diagonal pivots
$1/(2\nu),1/(2\nu),h/(2\nu^2)$, and its final scalar Schur
pivot is
\[
 4-2\nu j^2-\frac{\theta^2\nu^2}{2h}>0.
 \tag{FM22}
\]
Indeed $j^2\le1/4$, $0<\nu\le1$, $h\ge\nu$ and
$0\le\theta\le1$ in residue make the two subtracted terms
at most $1/2$ each. All entries before taking residues are finite
sums of bounded separated factors. The surviving strict
bounded-Gram lemma therefore supplies a finite tensor SOS.
Restoring the single square from FM20 and the whole positive
$\beta D_vN_x^2$ row gives the entire FM5. No mixed matrix inverse
or mixed square root has been used.

\paragraph{Testing, sign, and boundary joins.}
The argument includes $u=0$ and every real state without a state
normalization or a sign condition. Negative $u$ is changed to
$-u$ together with $(Y,Z)\mapsto(-Y,-Z)$; negative $v$ is
changed to $-v$ together with $(X,Z)\mapsto(-X,-Z)$.
FM5 is exactly invariant under these transformations. On FM1,
$e=\kappa h/\nu$ is infinitesimal and $C,A$ are unbounded.
The surviving SE10--SE12 theorem covers every unmatched
homogeneous Bob residue, including Bob zero and infinity,
and invokes the unchanged adjacent-$C,A$ theorem for
noninfinitesimal Alice tests. The only remaining Bob residue
has finite nonzero $v$ and $\bar v^2=\bar\ell/\bar\nu$,
which is the matched chart just proved. Thus the union covers
all projective tests on FM1. Real and imaginary state parts
are treated separately for complex states. All divisions and
normalizations used in these joins are pure source or pure
testing operations; all original source predicates remain.

\paragraph{Dependencies and limits.}
The analytical inputs are the surviving native source formulas
P1 converted as FM2--FM3, the full original parent, the strict
separated bounded-Gram lemma in \path{parent_reference.tex},
and \path{proofs/c144/small_e_unmatched_bob.tex}
with its preserved EB and adjacent-$C,A$ dependencies.
The exact budget is independently printed in
\path{c144_matched_bob_budget.tex}. The source endpoint
and unit-$\kappa$ sectors require separate complete arguments;
this theorem makes no claim about them.
